% Above: one control period on the bus, before and after the fix.
% Below: the life of a command.
\begin{tikzpicture}[font=\sffamily\scriptsize, x=1cm, y=1cm]

% ================================================================ the load
\node[chlabel, anchor=west] at (0,3.5) {one control period of 1.00 ms, four joints, drawn to scale from the arithmetic};

% the period boundary
\foreach \y in {2.0, 0.6}{ \draw[deadline] (4.8,\y-0.35) -- (4.8,\y+0.5); }
\node[tlabel, anchor=south, text=red!60!black] at (4.8,2.55) {the period ends here};

% row 1: everything at 1 kHz
\node[signame] at (-0.3,2.2) {command and state at 1 kHz};
\foreach \i in {0,...,3}{
  \node[taskrun, minimum width=10.7mm, anchor=south west, fill=hwcol] at ({\i*2.28},2.0) {cmd};
  \node[taskrun, minimum width=12.2mm, anchor=south west] at ({\i*2.28+1.07},2.0) {state};
}
\node[tlabel, anchor=west, text=red!60!black] at (9.4,2.2) {190 per cent: it does not fit};

% row 2: command at 250 Hz, arbitration at 1 Mbit
\node[signame] at (-0.3,0.8) {command at 250 Hz, arbitration at 1 Mbit/s};
\node[taskrun, minimum width=4.3mm, anchor=south west, fill=hwcol] at (0,0.6) {c};
\foreach \i in {0,...,3}{
  \node[taskrun, minimum width=7.6mm, anchor=south west] at ({0.5+\i*0.86},0.6) {state};
}
\node[tlabel, anchor=west] at (9.4,0.8) {38 per cent, and room to grow};

\node[umlnote, text width=50mm, anchor=north west] at (0,0.1)
  {The arbitration phase runs at the slow rate on every frame regardless of
   payload, so it is the expensive part. Raising the arbitration rate helps more
   than shrinking the layout, and doing both is what this volume would build for
   a real arm.};

% ================================================================ the command's life
\node[chlabel, anchor=west] at (0,-3.4) {and the life of one command};

\draw[taxis] (0,-4.8) -- (9.4,-4.8);
\node[taskrun, minimum width=26mm, anchor=south west] at (0.6,-4.6) {FRESH};
\node[taskrun, minimum width=26mm, anchor=south west, fill=hwcol] at (3.4,-4.6) {STALE};
\node[taskrun, minimum width=26mm, anchor=south west, fill=red!20] at (6.2,-4.6) {EXPIRED};
\node[tlabel, anchor=north, text width=26mm, align=center] at (1.9,-4.95) {the node acts on it};
\node[tlabel, anchor=north, text width=26mm, align=center] at (4.7,-4.95) {still valid, nothing newer};
\node[tlabel, anchor=north, text width=26mm, align=center] at (7.5,-4.95) {the named policy runs};

\foreach \x/\lab in {0.6/{received}, 3.4/{no newer frame arrives}, 6.2/{valid\_until passes}}{
  \draw[busstub] (\x,-4.05) -- (\x,-4.6);
  \node[tlabel, anchor=south, text width=26mm, align=center] at (\x,-4.0) {\lab};
}

\node[umlnote, text width=52mm, anchor=north west] at (0,-6.0)
  {A stale command is not an error: the master may simply have nothing new to
   say. An expired one is different, and the node does not merely stop acting on
   it, it enters the policy the command itself named, which by default is a
   controlled stop.};
\node[umlnote, text width=52mm, anchor=north west] at (5.6,-6.0)
  {The expiry time comes from the master, so the master chooses how long its own
   silence is tolerated. A node whose author guessed that number would be wrong
   for every system but the one it was written on.};
\end{tikzpicture}
